% ==========================
% REFERENCES
% ==========================

\chapter*{REFERENCES}
\renewcommand{\currentchapter}{REFERENCES}
\addcontentsline{toc}{chapter}{REFERENCES}

\noindent
Zhang, X., Li, Y. and Zhou, P. (2021) `MicroHECL: High-Efficient Root Cause Localization in Large-Scale Microservice Systems'. arXiv:2103.01782.

\vspace{1em}

\noindent
Zhao, H., Liu, J. and Wang, S. (2023) `Eadro: An End-to-End Troubleshooting Framework for Microservices on Multi-source Data'. arXiv:2302.05092.

\vspace{1em}

\noindent
Chen, R., Lin, Y. and Yu, M. (2023) `Towards Serverless Optimization with In-place Scaling'. arXiv:2311.09526.

\vspace{1em}

\noindent
Wang, K., Sun, J. and Li, T. (2024) `Root Cause Analysis in Microservice Using Neural Granger Causal Discovery (RUN)'. arXiv:2402.01140.

\vspace{1em}

\noindent
Liu, Z. and Tan, P. (2022) `Cushion: A Proactive Resource Provisioning Method to Mitigate SLO Violations for Containerized Microservices'. \textit{IET Communications}, 16(15), pp.~1802--1813. \url{https://doi.org/10.1049/cmu2.12464}.

\vspace{1em}

\noindent
Feng, Q., Zhao, Y. and Guo, M. (2024) `Root Cause Analysis for Microservices based on Causal Inference: How Far Are We?'. arXiv:2408.13729.

\vspace{1em}

\noindent
Huang, X., Li, D. and Zhao, F. (2024) `FaaSRCA: Full Lifecycle Root Cause Analysis for Serverless Applications'. arXiv:2412.02239.

\vspace{1em}

\noindent
Wu, J., He, L. and Zhang, Q. (2025) `LA-IMR: Latency-Aware, Predictive In-Memory Routing and Proactive Autoscaling for Tail-Latency-Sensitive Cloud Robotics'. arXiv:2505.07417.

\vspace{1em}

\noindent
Zhou, H., Chen, T. and Li, Y. (2025) `Simplifying Root Cause Analysis in Kubernetes with StateGraph and LLM (SynergyRCA)'. arXiv:2408.00803.

\vspace{1em}

\noindent
Patel, R., Kumar, A. and Singh, P. (2024) `A Comprehensive Survey on Root Cause Analysis in (Micro) Services: Methodologies, Challenges, and Trends'. arXiv:2407.10173.

\vspace{1em}

\noindent
Li, Y., Zhou, Z. and Hu, X. (2017) `A Trace-Based Performance Study of Autoscaling Workloads of Workflows in Datacenters'. arXiv:1711.08993.

\vspace{1em}

\noindent
Nguyen, T., Zhang, Y. and Wei, J. (2021) `Kubernetes Autoscaling: YoYo Attack Vulnerability and Mitigation'. arXiv:2105.00542.

\vspace{1em}

\noindent
Sun, L., Zhao, R. and Yang, J. (2024) `StatuScale: Status-aware and Elastic Scaling Strategy for Microservice Applications'. arXiv:2407.10173.

\vspace{1em}

\noindent
Rahman, M., Kim, D. and Park, S. (2021) `Reducing Cold Starts During Elastic Scaling of Containers in Kubernetes'. In: \textit{Proceedings of the 36th ACM Symposium on Applied Computing (SAC)}, pp.~420--427.

\vspace{1em}

\noindent
Tian, Y., Lee, K. and Zhou, M. (2024) `A Time Series-Based Approach to Elastic Kubernetes Scaling'. \textit{Electronics}, 13(3), p.~540.

\vspace{1em}

\noindent
Li, Z., Zhao, Y. and Wang, K. (2021) `Practical Root Cause Localization for Microservice Systems via Trace Analysis (TraceRCA)'. In: \textit{IEEE/ACM IWQoS 2021}. \url{https://doi.org/10.1109/IWQoS52092.2021.9521310}.

\vspace{1em}

\noindent
Masouros, D., Xydis, S. and Soudris, D. (2021) `Rusty: Runtime Interference-Aware Predictive Monitoring for Modern Multi-Tenant Systems'. \textit{IEEE Transactions on Parallel and Distributed Systems}, 32(12), pp.~3008--3020.

\vspace{1em}

\noindent
Xu, M., Zhao, D. and Chen, Y. (2022) `CoScal: Multi-faceted Scaling of Microservices with Reinforcement Learning'. \textit{IEEE Transactions on Network and Service Management}, 19(3), pp.~2201--2216. \url{https://doi.org/10.1109/TNSM.2022.3185938}.

\vspace{1em}

\noindent
Toka, L., Varga, P. and Marosi, A. (2021) `Machine Learning-Based Scaling Management for Kubernetes Edge Clusters'. \textit{IEEE Transactions on Network and Service Management}, 18(4), pp.~4567--4578. \url{https://doi.org/10.1109/TNSM.2021.3077278}.

\vspace{1em}

\noindent
Perez, J., Chen, L. and Gupta, S. (2018) `Investigating Performance Metrics for Scaling Microservices in CloudIoT Environments'. In: \textit{Proceedings of the ACM/SPEC International Conference on Performance Engineering (ICPE)}, pp.~233--240.